\documentclass[11pt]{article}
\usepackage{amsmath}
\usepackage{amssymb}
\usepackage{geometry}
\geometry{margin=1in}

\title{Arc Dipole Multipole Correction Procedure (corr\_MB\_v3.f)}
\author{Legacy code by S. Fartoukh (2010--2014)}
\date{}

\begin{document}
\maketitle

\section{Scope and Idea}
This document explains the legacy Fortran procedure \texttt{corr\_MB\_v3.f} (LHC Project Report 501 style) that generates arc corrector settings to compensate field errors of the main dipoles (MB/MBH). The code:
\begin{itemize}
  \item reads optics and error tables (\texttt{optics0\_MB.mad}, \texttt{MB.errors});
  \item converts integrated multipole errors into resonance-driving contributions using optics weights;
  \item solves for per-sector corrector strengths to cancel these contributions locally;
  \item distributes any remaining global skew terms using a pseudo-inverse;
  \item writes a MAD-X file \texttt{MB\_corr\_setting.mad}.
\end{itemize}
Eight arc sectors are used (periodicity of the LHC ring):
\[
\text{sectors } s=1\dots 8:\; (R1/L2,\ R2/L3,\ R3/L4,\ R4/L5,\ R5/L6,\ R6/L7,\ R7/L8,\ R8/L1).
\]
Beam direction is inferred from the order of sectors in \texttt{optics0\_MB.mad} and affects sign conventions.

\paragraph{Corrector families.} MQS (skew quads), MSS (skew sexts), MCS (normal sext spool), MCO (normal oct spool), MCD (normal deca spool), and MQT (normal quad trims, focus/defocus). Lengths are denoted $l.\mathrm{MQS}$, $l.\mathrm{MSS}$, etc. Integrated strengths printed in \texttt{MB\_corr\_setting.mad} are divided by these lengths.

\section{Optics Weighting}
For each MB slice the optics file provides $\beta_x,\beta_y, D_x, \Delta\mu_x, \Delta\mu_y$. Define
\[
\Delta\mu = 2\pi(\Delta\mu_x-\Delta\mu_y),\qquad c=\cos\Delta\mu,\qquad s=\sin\Delta\mu.
\]
Store per slice and sector $s$:
\begin{align*}
\text{b2 weights:}&\quad w_{b2,x}=\beta_x,\quad w_{b2,y}=\beta_y,\\
\text{a2 weights:}&\quad w_{a2,c}=\sqrt{\beta_x\beta_y}\,c,\quad w_{a2,s}=\sqrt{\beta_x\beta_y}\,s,\\
\text{a3 weights:}&\quad w_{a3,c}=\sqrt{\beta_x\beta_y}\,D_x\,c,\quad w_{a3,s}=\sqrt{\beta_x\beta_y}\,D_x\,s.
\end{align*}
Corrector responses are accumulated similarly:
\begin{itemize}
  \item MQS (skew quad): $a2_c(s) = \sum w_{a2}/(2\pi)\cdot \ell_{\mathrm{MQS}}$.
  \item MSS (skew sext): $a3_c(s) = \sum w_{a3}/(2\pi)$.
  \item MQT focus/defocus (normal quad trims): $b2_c^{F/D}(s) = \sum (\beta_x/2,\ -\beta_y/2)/(2\pi)\cdot \ell_{\mathrm{MQT}}$, with polarity chosen from sector parity, magnet index, and beam.
\end{itemize}
After counting slices, responses are renormalized (average over all corrector slices): $a2_c \gets a2_c\cdot 4/N_{\mathrm{MQS}}$, $b2_c^{F/D}\gets b2_c^{F/D}\cdot 8/N_{\mathrm{MQTF/D}}$; baseline strengths $k_{\mathrm{QTF/QTD}}$ are scaled accordingly.

\section{Error Aggregation}
From \texttt{MB.errors} each MB slice provides integrated multipoles $K_n, K_n^{(\mathrm{sk})}$. Per sector $s$ accumulate
\begin{align*}
b2_b &= \sum \bigl(w_{b2,x} K_1 - w_{b2,y} K_1\bigr)/(2\pi),\\
a2_b &= \sum w_{a2}\, K_1^{(\mathrm{sk})}/(2\pi),\\
a3_b &= \sum w_{a3}\, K_2^{(\mathrm{sk})}/(2\pi),\\
b3_b &= \sum K_2,\qquad b4_b = \sum K_3,\qquad b5_b = \sum K_4.
\end{align*}
Counts of MCS/MCO/MCD spools are stored to average later.

\section{First (Local) Correction}
\paragraph{Normal quad (b2).} Each sector uses the two MQT families (focus/defocus) to solve
\[
\begin{bmatrix}
b2_c^{F,x} & b2_c^{D,x}\\
b2_c^{F,y} & b2_c^{D,y}
\end{bmatrix}
\begin{bmatrix} k_F\\ k_D \end{bmatrix}
= - \begin{bmatrix} b2_b^x\\ b2_b^y \end{bmatrix},\qquad
(k_F,k_D)\to b2(s).
\]
\paragraph{Skew quad (a2).} There is one skew-quad knob per sector, but the error has two phase components (cosine/sine). Model it as a 2D vector equation
\[
a2_{\mathrm{loc}}(s)\, \mathbf{a2}_c(s) + \mathbf{a2}_b(s) \approx \mathbf{0},
\quad
\mathbf{a2}_c(s)=\bigl(a2_c^1(s),\,a2_c^2(s)\bigr),\;
\mathbf{a2}_b(s)=\bigl(a2_b^1(s),\,a2_b^2(s)\bigr).
\]
The best scalar $a2_{\mathrm{loc}}(s)$ that cancels along $\mathbf{a2}_c$ is the projection of $-\mathbf{a2}_b$ on $\mathbf{a2}_c$:
\[
a2_{\mathrm{loc}}(s) = -\frac{\mathbf{a2}_c\cdot\mathbf{a2}_b}{\|\mathbf{a2}_c\|^2}
= -\frac{\sum_i a2_c^i(s)\,a2_b^i(s)}{\sum_i (a2_c^i(s))^2}.
\]
Whatever is orthogonal to $\mathbf{a2}_c$ cannot be corrected locally; the remaining 2D error is accumulated as a residual for the global step:
\[
a2_{\mathrm{res}}^i \mathrel{+}= a2_{\mathrm{loc}}(s)\,a2_c^i(s)+a2_b^i(s).
\]
Net effect: kill the component parallel to the local response, carry the perpendicular component forward.
\paragraph{Skew sext (a3).} Same as a2 with $a3$ weights to get $a3_{\mathrm{loc}}(s)$ and residual vector $a3_{\mathrm{res}}$.
\paragraph{Normal sext/oct/deca.} Simple averages over spools:
\[
b3(s)=-\frac{b3_b(s)}{N_{\mathrm{MCS}}},\quad
b4(s)=-\frac{b4_b(s)}{N_{\mathrm{MCO}}},\quad
b5(s)=-\frac{b5_b(s)}{N_{\mathrm{MCD}}}.
\]
At collision energy, magnitudes are clipped to hardware limits $k^{\max}_{\mathrm{MCS}},k^{\max}_{\mathrm{MCO}},k^{\max}_{\mathrm{MCD}}$.

\section{Second (Global) Correction for a2 and a3}
After the local step, each sector still leaves a 2D residual vector (cos/sin) for skew quads and skew sexts:
\[
\mathbf{a2}_{\mathrm{res}}=\bigl(a2_{\mathrm{res}}^1, a2_{\mathrm{res}}^2\bigr),\qquad
\mathbf{a3}_{\mathrm{res}}=\bigl(a3_{\mathrm{res}}^1, a3_{\mathrm{res}}^2\bigr).
\]
Goal: distribute these two global components across all 8 sector knobs so that, when added up, they cancel the residual vector in a least-squares sense.

\paragraph{Response matrix.} Collect the sector responses into an $8\times 2$ matrix
\[
A = \begin{bmatrix}
a2_c^1(1) & a2_c^2(1)\\
\vdots & \vdots\\
a2_c^1(8) & a2_c^2(8)
\end{bmatrix}.
\]
Each row is the two phase components of the skew-quad response in one sector.

\paragraph{Pseudo-inverse.} The code builds
\[
B = -A^\top (A A^\top)^{-1}.
\]
This is the (negative) Moore–Penrose left pseudo-inverse of $A$; multiplying a 2D vector by $B$ yields an 8D vector of sector weights whose projection through $A$ best cancels that 2D vector in the least-squares sense.

Then the final per-sector strengths are
\[
a2(s) = a2_{\mathrm{loc}}(s) + \sum_{j=1}^2 B_{s j}\, a2_{\mathrm{res}}^j,\qquad
a3(s) = a3_{\mathrm{loc}}(s) + \sum_{j=1}^2 B_{s j}\, a3_{\mathrm{res}}^j.
\]
Interpretation:
\begin{itemize}
  \item $a2_{\mathrm{loc}}(s)$ (or $a3_{\mathrm{loc}}$) is what the sector did locally.
  \item $B_{sj} a2_{\mathrm{res}}^j$ is that sector's share of the global residual component $j$, chosen so that the weighted sum of all sectors reproduces (and cancels) the residual vector as closely as possible.
  \item Because there are 8 knobs but only 2 components to cancel, this distributes the correction smoothly across sectors rather than loading a single sector.
\end{itemize}


\section{Outputs}
\begin{itemize}
  \item \texttt{MB\_corr\_setting.mad} with:
\begin{itemize}
  \item Quadrupole trims: per-sector $\Delta$KQTF/KQTD and final KQTF/KQTD including baseline and beam sign.
  \item Skew quad (KQS), skew sext (KSS), normal sext (KCS), normal oct (KCO), normal deca (KCD) per sector, divided by element lengths.
\end{itemize}
  \item Sign conventions: depend on beam (B1 vs B2; a special ibeam=4 case if beam 2 and reversed order). Beam direction flag $bv$ flips certain families; the sign choice matches the MAD-X convention for the two beams.
  \item Saturation: at collision energy, spools are clipped to hardware limits ($k^{\max}_{\mathrm{MCS/MCO/MCD}}$) before printing.
\end{itemize}

\section{Reference Run}
Reference inputs are \texttt{temp/optics0\_MB.mad} and \texttt{temp/MB.errors}; the legacy output is \texttt{temp/MB\_corr\_setting.mad} in the same folder, matching the method above and the strategy described in LHC Project Report 501.

\paragraph{Reading tip.} If you want to trace a specific term: (i) find its weight in the optics section, (ii) see how it is multiplied by the corresponding $K_n$ column in \texttt{MB.errors}, (iii) note which corrector family and sector solve it (local) or receive a share via the pseudo-inverse (global skew terms).

\end{document}
